\documentclass[11pt]{article}
\usepackage{manual_docs_style}
\hypersetup{hidelinks}
\externaldocument{build/generate_question_stage}[generate_question_stage.pdf]
\externaldocument{build/agentic_rollout}[agentic_rollout.pdf]

\begin{document}

\docTitle{Website}
\phantomsection\label{docstart:website}
\label{sec:website}


\phantomsection\label{sec:website-document-contents}
\section*{Document Contents}
This document defines the public-facing design and the technical contract of the TraceBench website.
It also describes the \hyperref[sec:website-deployment]{production deployment flow} needed to publish the site and benchmark artifacts.
Benchmark data, public results, and accepted submission metadata are stored on the Hugging Face Hub and rendered by a static single-page application hosted through GitHub Pages. Trajectory files are part of the submission schema, but must only be advertised as downloadable after they have been published in the corresponding archive.
\newcommand{\websiteContentsEntry}[2]{\hyperref[#1]{#2\hfill\pageref*{#1}}}
\begin{center}
\renewcommand{\arraystretch}{1.18}
\begin{tabular}{|P{0.43\textwidth}|P{0.43\textwidth}|}
\hline
\textbf{\websiteContentsEntry{sec:website-design-brief}{Design Brief}} &
\textbf{\websiteContentsEntry{sec:website-technical-brief}{Technical Brief}} \\
\hline
\websiteContentsEntry{sec:website-overview}{Overview} &
\websiteContentsEntry{sec:website-hosting-structure}{GitHub Pages structure} \\
\websiteContentsEntry{sec:website-user-admin-actions}{User and Admin actions} &
\websiteContentsEntry{sec:website-raw-submission-bundle}{Raw submission bundle schema} \\
\websiteContentsEntry{sec:website-global-style}{Global style} &
\websiteContentsEntry{sec:website-routing}{Website structure and routing} \\
\websiteContentsEntry{sec:website-page-overview}{Overview} &
\websiteContentsEntry{sec:website-deployment}{Deployment workflow} \\
\websiteContentsEntry{sec:website-page-results}{Results} &
\websiteContentsEntry{sec:website-public-data-schema-appendix}{Appendix: public/data schema} \\
\websiteContentsEntry{sec:website-page-environments}{Environments} &
\websiteContentsEntry{sec:website-page-get-started}{Get Started and Submit} \\
\websiteContentsEntry{sec:website-page-contributors}{Contributors, Citation, and Contact} & \\
\hline
\end{tabular}
\end{center}

\phantomsection\label{sec:website-design-brief}
\section*{Design Brief}
\phantomsection\label{sec:website-overview}
\subsection*{Overview}
The goals of the website are to:
\begin{itemize}
    \item Make the benchmark understandable before a reader opens the paper or the repository.
    Specifically, explain what TraceBench measures, how tasks are generated, what results were observed, how users can inspect examples, and how external evaluations can be submitted.
    \item Enable people to download results and datasets, and to inspect trajectories when the released archive contains them.
    \item Allow third parties to submit their agentic scores.
\end{itemize}
The website is hosted at \code{tracebench.github.io}.
Results, complete result tables, benchmark data, and downloadable artifacts are hosted on the Hugging Face Hub. The website must not imply that raw trajectories are available for a submission unless its public archive includes the trajectory files (or explicit trajectory URLs).
The website should be both desktop- and mobile-friendly.



\phantomsection\label{sec:website-user-admin-actions}
\subsection*{User and Admin actions}
The user actions are
\begin{itemize}
\item Inspect the leaderboard.
\item Download complete result tables and benchmark datasets stored on Hugging Face; download trajectories only for submissions that publish them.
\item Submit new results. 
This works as follows:
\begin{itemize}
  \item The submitter opens a GitHub pull request with a structured submission with schema \hyperref[sec:website-raw-submission-bundle]{Raw submission bundle schema}.
  \item A maintainer-side submission validation workflow validates the submission format and coverage, reads the submitted scores and trajectories, and writes a review artifact containing validation reports, proposed canonical result rows, trajectory indexes, artifact indexes, and a proposed submission archive.
  A maintainer then reviews this artifact as described in \hyperref[sec:website-admin-review-submissions]{Inspect and accept/reject submissions}.

\end{itemize}
\end{itemize}
The admin actions are
\begin{itemize}
  \item Populate the Hugging Face repository \code{eth-siplab/tracebench} with the initial Hugging Face dataset state from the current TraceBench repository state.
  \item Regenerate the website JSON from Hugging Face with GitHub Actions  
  \item \phantomsection\label{sec:website-admin-review-submissions}Inspect and accept/reject submissions: Maintainers access the submission artifact from the pull request's GitHub Actions run and inspect the validation report, proposed canonical result rows, trajectory index, artifact index, and proposed accepted submission archive.
  If the submission is approved, a maintainer merges the pull request and triggers a separate maintainer-only GitHub Actions workflow for publication.
  This publication workflow \code{publish-submission.yml} does:
  \begin{itemize}
    \item using an authenticated Hugging Face token, rebuilds the accepted submission archive from the merged repository state, appends the accepted canonical result rows to \texttt{results.parquet}, validates the updated dataset, and uploads \texttt{submissions/<submission-id>.zip} and the updated \texttt{results.parquet} to Hugging Face.
    \item regenerates the browser JSON from Hugging Face. 
  \end{itemize}
  \item Add new \termSimulators by uploading their question bundles and metadata to Hugging Face, then regenerating the website JSON from that canonical dataset state.
  This creates a new version of the public simulator catalogue.
  \item Make website content/style changes 
\end{itemize}


\phantomsection\label{sec:website-global-style}
\subsection*{Global style}

The intended style is minimalistic and simple. 
It should follow the design principle of `Perfection is not when there is nothing more to add, but there is nothing left to take away'.
White background, simple navigation, strong typography, rectangular figures, thin divider lines, and a blue link style consistent with ETH SIPLAB pages should be the standard.

The website should use a restrained research-lab style:

\begin{itemize}
    \item \textbf{Layout.} Use a full-width main content area with a centered maximum-width content column; figures, plots, and interactive examples should span the full available content width and scale cleanly on mobile.
    \item \textbf{Typography.} Use large bold section titles, short descriptive subtitles, and compact body text.
    \item \textbf{Navigation.} Use a simple top header menu with text links. The website must be implemented as a single-page application with client-side routing. These links should update the browser URL and render the corresponding route without loading separate HTML documents or merely scrolling to sections within one long page.
    \item \textbf{Visuals.} Prefer rectangular plots and tables.
    \item \textbf{Color.} Use mostly black, white, and grey. Use blue for links and active menu items. Use orange only as a secondary accent.
    \item \textbf{Components.} Prefer thin rules, tables, and flat panels over heavy cards or shadows.
    \item \textbf{Section labels.} Do not use eyebrow labels or small uppercase pre-headings above page titles or section titles. 
    Each page should begin directly with its main title and short subtitle/description.
    Section hierarchy should be communicated through headings, spacing, and thin divider lines rather than additional decorative labels.
\end{itemize}

\subsubsection*{Header}

The top navigation should contain the TraceBench wordmark on the left and a horizontal menu:

\begin{center}
\texttt{Home \quad Results \quad Environments \quad Get Started \quad Contributors \quad Paper (arXiv) \quad Code \quad Hugging Face}
\end{center}

The active page should be highlighted in blue.
The header should have a thin horizontal rule below it, not a large navigation box.
Paper, Code, and Hugging Face are external links in the same header.
The Hugging Face link should point to \code{hf_dataset_url} from \code{site.json}.



\section*{Page specifications}
\phantomsection\label{sec:website-page-overview}
\subsection*{Page 1: Overview}
The homepage should immediately explain TraceBench, show the key benchmark statistics, provide a compact example task visual, and end with the key features.

\subsubsection*{Hero}

The hero should be a single full-width introductory block within the main centered content column.
It should not use a two-column layout and should not place secondary panels, tables, statistics, benchmark-scope content, task-definition content, or design-principle content beside the title.
The hero should contain only the main title and the two-paragraph subtitle/description.
All other homepage content should appear below the hero.

The hero should contain the title:

\begin{quote}
\agentProposedEdit{agent-remove-4f28c3a4-26c6-5a0c-9c4a-9a33e7a3a37c}{superseded title}{\textbf{TraceBench: A Controlled Evaluation for Agentic Root-Cause Analysis of Time Series}}{}
\end{quote}

The subtitle should explain the benchmark in two short paragraphs:

\begin{quote}
A benchmark for evaluating whether LLM agents can perform evidence-grounded analysis of multivariate time series.

Tasks are generated from physical simulators with known interventions.
Agents must combine textual descriptions, noisy observations, and optional labeled examples to identify what changed, or that no intervention occurred.
\end{quote}


\subsubsection*{Homepage content order}

Below the hero, the landing page should present the homepage sections in the following order:

\begin{enumerate}
  \item Illustrative Example (GIF and Example)
  \item Statistic row
  \item Key features
\end{enumerate}

The \texttt{Key features} section should appear near the bottom of the landing page, below the example task visual.


\subsubsection*{Illustrative example}
This homepage section explains to new users how samples are generated and how agents are prompted during evaluation on the generated sample.
The section should contain two parts: \textbf{Data Simulation}, showing the generation process as a GIF, and \textbf{Agentic Question}, showing an example time-series plot and the corresponding agent-facing question.

\paragraph{Data Simulation}
The data simulation part comprises GIF of the BallDrop simulator as the first content section below the hero and above the statistic row.
The GIF should be displayed directly on the white page background, without an additional outer border, frame, card, or enclosing rectangle around the image.
On the homepage, the GIF should be displayed at the full width of the main content column, matching the text column width.
The GIF should be stored and served in the website repository at \code{assets/media/ball-drop-bounce.gif}.

The design brief used to generate the GIF is located \hyperref[sec:website-ball-drop-gif-prompt]{here}.

\paragraph{Question}
The example time-series plot should be rendered from the generated environment JSON files under \code{public/data/environments/BallDrop/description.json} and \code{public/data/environments/data_main_page.json}.
The plot should be rendered like the \hyperref[sec:website-environment-specific-page]{environment-specific examples}, with two differences: (i) zoom should be disabled and (ii) it must not have any dashed intervention line.
The layout should be single-column: plots specific control above the plot, the plot across the full content width, prompt specific control above the prompt,  the prompt panel below.

Controls should be grouped into labelled rows:
Plots specific control:
\begin{DocCode}
  Noise: [NO NOISE] [LOW NOISE] [HIGH NOISE]
\end{DocCode}
The controls should update the plot when switching between \texttt{No Noise}, \texttt{Low Noise}, and \texttt{High Noise}.

Prompt specific control:
\begin{verbatim}
Task: [CODE] [DIRECT]
Context: [NO CONTEXT] [HIGH CONTEXT]
Examples: [NO EXAMPLES] [ONE EXAMPLE] [MULTIPLE EXAMPLES]
\end{verbatim}
The \texttt{Task}, \texttt{Context}, and \texttt{Examples} controls should update the displayed prompt.

The highlighted controls showing the active state.

Together, these interactions should be just enough to demonstrate task mode and condition-axis changes.
There should be two buttons show "Show Prompt" and "Reveal Answer".
The prompt specific control and the prompt should appear when "Show Prompt" is clicked.
The answer should pop up. if "Reveal Answer" is clicked.
The prompt visualized is defined in \hyperref[sec:website-prompt-appendix]{Appendix: Website Prompt Examples}.

\subsubsection*{Statistic row}

The homepage should include two compact statistics in the following exact order:

\begin{center}
\begin{tabular}{ll}
\hline
\textbf{Statistic} & \textbf{Meaning} \\
\hline
\textbf{3 physical simulators} & BallDrop, BounceBall, MassSlide \\
\textbf{4 experimental conditions} & Noise, Task output format, Context, Number of training examples \\
\hline
\end{tabular}
\end{center}



\paragraph{Key features}
\begin{itemize}
  \item \textbf{Physics-grounded generation.} Samples are produced by simulators rather than by static datasets, giving reliable intervention labels and controllable task variants.
  \item \textbf{Evidence-grounded agent evaluation.} Correct answers require inspecting the observed signals; textual context alone is not sufficient.
  \item \textbf{Controlled difficulty axes.} TraceBench varies textual context, observation noise, labeled examples, and answer interface in matched experimental conditions.
\end{itemize}



\phantomsection\label{sec:website-page-results}
\subsection*{Page 2: Results}
The Results page should start with one adjustable leaderboard where different filters can be selected. 
By default it should open \textbf{Code, Low Noise, High Context, Three Examples}, which we consider the canonical result.
The other experimental conditions can be accessed using filters.
A submission is accepted and thus visualized in the a leaderboard view only if it includes exactly five completed seeds for each simulator under the selected condition. 
With the current benchmark setup of three simulators, this means a submission must include 15 valid runs for that condition before it is displayed.


\subsubsection*{Table behavior}
The Results page should display a leaderboard table for one selected experimental condition.
By default, the selected condition should be Code, Low Noise, High Context, and Three Examples.

The table has six columns: rank, agent, model, score, date, and submitter.
Each column supports sorting.

Above the table, the user can select the displayed experimental condition using condition selectors for task mode, noise level, context level, and number of examples.
These selectors determine which leaderboard scores are shown.

In the top-right area above the table, there is an additional button: Compare.
When Compare is enabled, the score column should display both scores for each row, one for the primary selected condition and one for the comparison condition.
The score display should make the direction of comparison clear, for example as primary score \(\rightarrow\) comparison score, or as two labeled values.

\subsubsection*{Clickable leaderboard rows}
Each leaderboard row should be clickable and should navigate to a result-detail route within the single-page application. 
The route should preserve the active leaderboard filter, for example:

\begin{DocCode}
/results/<submission-id>/?task_mode=Code&noise=Low&context=High&examples=Three
\end{DocCode}

By default, the result-detail page should display the same condition that was active when the user clicked the leaderboard row.
The page should also allow the user to change the task mode, noise level, context level, and number of examples without leaving the result-detail route.

The result-detail page should show the selected agent, date, submitter, displayed score, active filter scope, per-condition results, per-seed results, trajectory links, and downloadable artifact links. 
A displayed score should appear only when the selected condition contains exactly five valid seeds for each required simulator.


The browser URL should update when a user opens a result-detail page.
The route should remain inside the single-page application and should not trigger a full-page reload.


\phantomsection\label{sec:website-page-environments}

\subsection*{Page 3: Environments}
The Environments page should present all simulators in a square-card grid. 
Each card should show the simulator name and the one-line description from \code{description.json::short_one_line_description}. 
Clicking a card should navigate within the single-page application to an environment-specific route, for example \code{/environments/BallDrop/}.

\phantomsection\label{sec:website-environment-specific-page}
\paragraph{Environment-specific page}
Each environment-specific page explains one simulator in detail. 
The page should begin with a full-width Plotly time-series plot, followed by the agent prompt for the selected task condition. 
As on the homepage example visual, the user should be able to choose between task mode and condition-axis settings. 
However, this page should:
\begin{itemize}
  \item Include a dropdown for selecting different samples from the generated environment JSON files. 
  \item \agentProposedEdit{agent-remove-373beebd-5c9b-54a8-ab46-5342c7dc9441}{superseded by the operational legend-toggle contract}{Allow to choose between multiple signals (and not just one), with plot axis rescaled accordingly.}{}
  \item \agentProposedEdit
    {agent-add-8fcee341-2756-54ad-b411-ad8027fa543d}
    {Add this agent-authored block for review.}
    {}
    {Environment plots allow multiple signals to be shown or hidden through the Plotly legend.}
  \item \agentProposedEdit
    {agent-add-4ac30cdd-c1ca-57cf-aeee-d977482b3647}
    {Add this agent-authored block for review.}
    {}
    {After a single-click visibility toggle or double-click trace isolation, the plot recomputes the y-axis range from the currently visible signals.}
  \item \agentProposedEdit
    {agent-add-ba2b250c-0e68-5610-ba5e-6c694105bac8}
    {Add this agent-authored block for review.}
    {}
    {Manual y-axis zoom and pan remain in effect until the next legend visibility change.}
  \item The prompt should be the correct one 
\end{itemize}
\agentProposedEdit
  {agent-add-621a1db1-3653-5581-9823-983cce191c0d}
  {Add this agent-authored block for review.}
  {}
  {The y-axis autorange behavior is implemented in \code{assets/app.js} at website commit \code{41df7900e4141b2d66ffcd69cd6d39aa84dca336} and was verified on the deployed \code{/environments/BallDrop/} page.}
The sampled are named 1, 2, 3, 4, 5 and are mapped to \code{data_<n>} where n is the number.
Changing the controls should update the plotted sample and the displayed prompt without leaving the route.
The data-loading contract for environment-specific pages is specified in the \hyperref[sec:website-generated-json]{Generated JSON files} section of the Technical Brief.

\phantomsection\label{sec:website-page-get-started}
\subsection*{Page 4: Get Started and Submit}

The get-started page should be practical.
It should contain a large copy-pasteable code block and a clear submission flow.

\subsubsection*{Run locally}

The quick-start section should contain:

\begin{verbatim}
# Clone and install
git clone <TRACEBENCH_PUBLIC_REPOSITORY_URL>
cd TraceBench
python -m venv env
source env/bin/activate
pip install -e .

# Launch local explorer to explore the data locally
bash web_model_explorer/start.sh

# Run one agentic evaluation
python workflows/rollout/question_run_orchestrator.py \
  --tasks-dir TraceBench_questions \
  --model BallDrop \
  --agent-id <AGENT_ID> \
  --row-slug <ROW_SLUG>
\end{verbatim}

The code block should be visually prominent because this page is primarily a usage page.

The page should also show how to submit a new result.

\begin{quote}
To submit a new result, open a GitHub pull request with one structured submission bundle.
The bundle should include submission metadata, one folder for each submitted question/run, the corresponding \code{scores.json} and \code{atif_trajectory.json} files, and the artifacts produced during the run.
The exact file layout is defined in the \hyperref[sec:website-raw-submission-bundle]{Raw submission bundle schema}.
\end{quote}


\phantomsection\label{sec:website-page-contributors}
\subsection*{Page 5: Contributors, Citation, and Contact}

The final page should make the website feel complete, citable, and easy to contact.
It should include contributors, citation information, and contact links.

\subsubsection*{Contributors}

The page should include a clear \textbf{Contributors} section listing the project contributors, affiliations, and optionally equal-contribution notes.

\subsubsection*{Citation}

The citation should be shown in a dark code block

\subsubsection*{Contact}

The final page should include a clear \textbf{Contact} section.
This was added based on the print annotation.
\begin{center}
\texttt{Email \quad Repository \quad Affiliation}
\end{center}

All public site metadata for this page is stored in \code{site.json} in the GitHub Pages repository.

\phantomsection\label{sec:website-technical-brief}
\section*{Technical Brief}

\subsection*{Target destinations}
The target destinations are:
\begin{itemize}
\item Website repository: https://github.com/tracebench/tracebench.github.io
\item Deployed website: https://tracebench.github.io/
\item Code repository: https://github.com/TommasoBendinelli/TraceBench.git
\item Hugging Face dataset repository: https://huggingface.co/datasets/eth-siplab/tracebench
\end{itemize}
The local website checkout for implementation should be a sibling repository \code{/Users/tbe/repos/tracebench.github.io}.

\newcommand{\websiteTreeSummaryJson}{\hyperref[sec:website-summary-json]{summary.json}}
\newcommand{\websiteTreeLeaderboardJson}{\hyperref[sec:website-leaderboard-json]{leaderboard.json}}
\newcommand{\websiteTreeEnvironmentDescriptionJson}{\hyperref[sec:website-environment-description-json]{description.json}}
\newcommand{\websiteTreeEnvironmentDataJson}{\hyperref[sec:website-environment-data-json]{\textless data\_1\textgreater.json}}
\newcommand{\websiteTreeGenerateWebsiteData}{sync\_hf\_to\_website\_data.py}
\newcommand{\websiteTreeDeployWorkflow}{\hyperref[sec:website-deployment]{deploy.yml}}


\subsection*{Schema and Structure}
\phantomsection\label{sec:website-hosting-structure}
\subsubsection*{GitHub Pages structure}
\begin{LinkedDocCode}
tracebench.github.io/
  index.html
  404.html
  site.json
  assets/
    app.js
    styles.css
    media/
      ball-drop-bounce.gif
  public/
    data/
      \websiteTreeSummaryJson
      \websiteTreeLeaderboardJson
      environments/
        <environment_1>/
          \websiteTreeEnvironmentDescriptionJson
          \websiteTreeEnvironmentDataJson
          ...
      ...
  scripts/
    \websiteTreeGenerateWebsiteData
  .github/
    workflows/
      \websiteTreeDeployWorkflow
\end{LinkedDocCode}
Note that:
\begin{itemize}
  \item  All public site metadata is stored in \code{site.json} at the root of the GitHub Pages repository.
  The file is committed with the website source and is not generated from Hugging Face.
  \begin{DocCode}
  {
    "website_url": "https://tracebench.github.io/",
    "code_url": "https://github.com/TommasoBendinelli/TraceBench",
    "hf_repo": "eth-siplab/tracebench",
    "hf_dataset_url": "https://huggingface.co/datasets/eth-siplab/tracebench",
    "issues_url": "https://github.com/TommasoBendinelli/TraceBench/issues",
    "affiliation_url": "https://siplab.org/",
    "paper_url": "https://arxiv.org",
    "contributors": [
      {
        "name": "Tommaso Bendinelli",
        "affiliation": "ETH Zurich SIPLAB / CSEM",
        "url": ""
      }
    ],
    "contact": {
      "label": "GitHub issues",
      "url": "https://github.com/TommasoBendinelli/TraceBench/issues"
    },
    "citation": "@inproceedings{tracebench2026,\n  title={TraceBench: A Controlled Evaluation for Agentic Root-Cause Analysis of Time Series},\n  author={Tommaso Bendinelli and Artur Dox and Christian Holz},\n  year={2026}\n}"
  }
  \end{DocCode}

  \item Files in \code{public/data/*} are generated by \code{scripts/sync_hf_to_website_data.py} which is triggered by the GitHub Action \code{deploy.yml}.
  The detailed schema examples for these files are collected in \hyperref[sec:website-public-data-schema-appendix]{Appendix: public/data schema}.
  \item Each environment-specific page loads its metadata from \code{public/data/environments/<environment>/description.json} and its plotted samples from \code{public/data/environments/<environment>/data_<n>.json}.
\end{itemize}




\phantomsection\label{sec:website-raw-submission-bundle}
\subsubsection*{Raw submission bundle schema}

A submitted result should be a structured bundle committed through a GitHub pull request.
Each submission bundle is scoped to one agent/submission at a time to keep validation and publication atomic.
Each \code{<question_id>} folder represents one submitted seed/run.

\begin{DocCode}
<submission_id>/
  seed_coverage.json
  <question_id>/
    scores.json
    atif_trajectory.json
    artifact/*
  <question_id>/
    scores.json
    atif_trajectory.json
    artifact/
      ...
\end{DocCode}
These files are the result of running the rollout as defined in \hyperref[sec:agentic-rollout-materialized-output]{Agentic Rollout: Materialized Output}.


\subsection*{Hugging Face structure}
\begin{DocCode}
eth-siplab/tracebench/
  README.md
  LICENSE
  questions/
  results.parquet
  submissions/
    <submission-id>.zip
\end{DocCode}
\begin{itemize}
    \item \code{README.md}: Dataset card for the Hugging Face repository. 
    It explains what TraceBench is, how the benchmark questions are organized, how to load/use the files, how submissions are stored, and how to cite the benchmark.
    The submission instructions should use the same public-facing message shown on the \hyperref[sec:website-page-get-started]{Get Started and Submit} page.
    \item \code{LICENSE}: License for the released benchmark data, question bundles, public results, and submission artifacts.
    \item \code{questions/}: Public TraceBench question bundles.
    See \hyperref[docstart:generate_question_stage]{Stage: Generate Question} for the exported bundle schema and generation workflow.
    \item \code{results.parquet}: Canonical public results table. 
    It contains one row per evaluated public run, usually identified by submission, agent/model, simulator, task condition, and seed. 
    The website leaderboard JSON is generated from this file by grouping and filtering the rows, for example requiring five completed seeds for each simulator before showing a public leaderboard score.
    \item \code{submissions/}: Archive of accepted public submissions. 
    Each submission is stored as a zip file. 
    The schema of each of these is defined as \hyperref[sec:website-raw-submission-bundle]{Raw submission bundle schema}.
    The zipping happens after validation.
\end{itemize}

\section{Scripts}
\subsection*{\code{validate_submission.py}}
This script validates the submission format and required seed coverage, reads the submitted scores and trajectory files, and produces a GitHub Actions review artifact containing the information needed for maintainer approval or rejection.

\subsection*{\code{scripts/sync_hf_to_website_data.py}}

\code{scripts/sync_hf_to_website_data.py} converts the canonical Hugging Face benchmark state into the small browser-ready JSON files consumed by the website.
It should read only from Hugging Face: the public results table, generated question metadata, and accepted submission archives.
It then writes only the compact public data needed by the static site.
During this sync, it renders the website prompt examples from \code{questions/<environment>/questions.json} by calling \code{shared/prompts.py::render_tracebench_agent_prompt(...)} from the public TraceBench repository checkout \code{https://github.com/TommasoBendinelli/TraceBench.git}.

The script writes the files listed in \hyperref[sec:website-generated-json-file-list]{Generated JSON files}.

\subsection*{\code{scripts/generate_base_results.py }}
Takes as input a command \code{--output} to specify the folder where it's going to be created
It's a one-time or maintainer-triggered publication script used to create the initial Hugging Face dataset state from the current TraceBench repository state.
This script belongs to the TraceBench repository or to a controlled maintainer publication environment.
It is not part of the normal GitHub Pages deployment workflow
\subsection*{\code{serve.py}}
To serve the website locally

\section{GitHub Actions}
\subsection*{\code{validate-submission.yml}}
   Trigger: pull request with a submitted result bundle
   Purpose: validate the submission and create the review artifact
\subsection*{\code{publish-submission.yml}}
  Trigger: manual maintainer-only \code{workflow\_dispatch} after PR approval/merge
  Purpose: publish accepted result artifacts to Hugging Face and triggers deploy.yml afterward
\subsection*{\code{deploy.yml}}
  Trigger: push to website main, manual \code{workflow\_dispatch}, or triggered by publish-submission.yml
  Purpose: regenerate public/data/*.json from Hugging Face and deploy GitHub Pages

\phantomsection\label{sec:website-routing}
\section*{Website structure and routing}

The website must be implemented as a single-page application (SPA), not as a collection of separate static HTML pages and not as a single long scrolling page with anchor sections.
Each top-level navigation item should map to a client-side route handled by the SPA router:

\begin{itemize}
    \item \texttt{/}: \hyperref[sec:website-page-overview]{Overview}
    \item \texttt{/results/}: \hyperref[sec:website-page-results]{Results}
    \item \texttt{/environments/}: \hyperref[sec:website-page-environments]{Environments}
    \item \texttt{/get-started/}: \hyperref[sec:website-page-get-started]{Get Started and Submit}
    \item \texttt{/contributors/}: \hyperref[sec:website-page-contributors]{Contributors, Citation, and Contact}
\end{itemize}

The SPA should also support result-detail routes for leaderboard records:

\begin{itemize}
    \item \texttt{/results/}: Results page with one filterable leaderboard
    \item \texttt{/results/\textless submission-id\textgreater/}: detail page for one displayed leaderboard record
\end{itemize}

The navigation menu should link to these routes directly.
Route changes should preserve the SPA shell and should not trigger full-page reloads.

\subsection*{SPA deep links}

\subsubsection*{Background of the problem}
Because the website is a single-page application, client-side routes such as these must work when opened directly or refreshed in the browser and not just when we get there through the app. 
Open directly is when the user types or pastes a nested URL into the browser address bar.
Refreshed means the user is already the correct link, and then it clicks the refresh button.
\agentProposedEdit{agent-remove-64d47c85-8a92-508c-b1fc-4f8344ab9288}{deep-link route files are now materialized}{Specifically the website has only:}{}
\begin{DocCode}
  index.html
  404.html
  assets/app.js
  assets/styles.css
\end{DocCode}
So without proper handling, when somebody opens a URL such as \code{https://tracebench.github.io/results/gpt-55-agent-direct-20260501/}, GitHub Pages looks for a matching physical file.
Since that file does not exist, GitHub Pages serves \code{404.html}.


\subsection*{Publishing the website}

From a local checkout of the website mockup:

\begin{verbatim}
cd website_local_version

git init
git branch -M main
git remote add origin git@github.com:tracebench/tracebench.github.io.git

git add .
git commit -m "Initial TraceBench website"
git push -u origin main
\end{verbatim}

In GitHub, configure Pages to deploy through GitHub Actions:

\begin{verbatim}
Repository -> Settings -> Pages -> Build and deployment -> Source -> GitHub Actions
\end{verbatim}




\subsection*{Deployment}
\phantomsection\label{sec:website-deployment}


\subsection*{Required secrets}
  The following GitHub Actions secrets are used by the publication and deployment workflows.
  \begin{itemize}
    \item \code{HF_TOKEN}: Hugging Face token with access to \code{eth-siplab/tracebench}.
    It is required in the TraceBench repository for \code{publish-base-dataset.yml} and \code{publish-submission.yml}, where it must have write access.
    It is also required in the \code{tracebench.github.io} repository if \code{deploy.yml} reads from a private Hugging Face dataset.
    If the Hugging Face dataset is public, \code{deploy.yml} may read without \code{HF_TOKEN}.
    \item \code{TRACEBENCH_GITHUB_TOKEN}: Optional GitHub token used only if maintainer publication workflows should automatically trigger \code{deploy.yml} in the separate \code{tracebench.github.io} repository.
    If maintainers trigger the website deployment manually, this secret is not required.
  \end{itemize}

\subsubsection*{Phase 1: publish the canonical Hugging Face dataset}
Prerequisite: a local public question-bundle directory containing \code{BallDrop/}, \code{BounceBall/}, and \code{MassSlide/}, with each bundle matching the schema of \hyperref[docstart:generate_question_stage]{Stage: Generate Question}.

From the TraceBench repository or controlled maintainer environment:
\begin{DocCode}
cd /Users/tbe/repos/TraceBench
python scripts/generate_base_results.py \
    --question-bundles-dir /path/to/public/question/bundles \
    --output /tmp/tracebench_hf_dataset
\end{DocCode}

Create the Hugging Face repository \code{eth-siplab/tracebench}.
Create an \code{HF_TOKEN} with read and write access to this dataset repository.
Upload the generated folder in \code{/tmp/tracebench_hf_dataset} to \code{eth-siplab/tracebench} using Hugging Face tooling.

If the upload is performed through GitHub Actions, configure \code{HF_TOKEN} as a repository secret in the TraceBench GitHub repository.

\subsubsection*{Phase 2: prepare the website repository}
Then trigger:
\begin{DocCode}
cd /Users/tbe/repos/tracebench.github.io
python scripts/sync_hf_to_website_data.py --hf-repo eth-siplab/tracebench
\end{DocCode}

If \code{eth-siplab/tracebench} is private, configure \code{HF_TOKEN} as a repository secret in the \code{tracebench.github.io} GitHub repository so \code{deploy.yml} can read from Hugging Face.
If the Hugging Face dataset is public, \code{deploy.yml} may read the website data without \code{HF_TOKEN}.

Initialize and push the website repository if it is not already cloned:
\begin{DocCode}
git init
git branch -M main
git remote add origin git@github.com:tracebench/tracebench.github.io.git
git add .
git commit -m "Initial TraceBench website"
git push -u origin main
\end{DocCode}

\phantomsection
\section*{Appendix: public/data schema}
\label{sec:website-generated-json}\label{sec:website-generated-json-file-list}\label{sec:website-public-data-schema-appendix}

In the following we describe the structure of each file in \code{public/data/*}.
\phantomsection\label{sec:website-summary-json}
\begin{DocCode}
public/data/summary.json
{
  "num_simulators": 3,
  "simulators": ["BallDrop", "BounceBall", "MassSlide"],
  "task_modes": ["Code", "Direct"],
  "axes": ["noise", "context", "examples"],
  "canonical_scope": {
    "task_mode": "Code",
    "noise": "Low",
    "context": "High",
    "examples": "Three Examples",
    "required_seeds_per_simulator": 5
  }
}
\end{DocCode}
\phantomsection\label{sec:website-leaderboard-json}
\begin{DocCode}
public/data/leaderboard.json
{
  "generated_at": "2026-05-31T00:00:00Z",
  "rows": [
    {
      "rank": 1,
      "submission_id": "agent-x-2026-05-01",
      "agent": "Agent X",
      "model": "Model Y",
      "score": 0.82,
      "date": "2026-05-01",
      "submitter": "Lab Z",
      "scope": {
        "task_mode": "Code",
        "noise": "Low",
        "context": "High",
        "examples": "Three Examples"
      }
    }
  ]
}
\end{DocCode}
\phantomsection\label{sec:website-environment-description-json}
\begin{DocCode}
public/data/environments/BallDrop/description.json
{
  "name": "BallDrop", # comes from experiment_config.json::paper_facing_name, must match public/data/environments/<this>/description.json
  "short_one_line_description": "A ball is dropped from a height and observed over time.",
  "prompt_combinations": [
    {
      "desc_level": "high",
      "task_type": "direct",
      "training_samples": "none",
      "agent_instruction": "..."
    },
    {
      "desc_level": "high",
      "task_type": "code",
      "training_samples": "none",
      "agent_instruction": "..."
    },
    {
      "desc_level": "high",
      "task_type": "direct",
      "training_samples": "one",
      "agent_instruction": "..."
    },
    {
      "desc_level": "high",
      "task_type": "code",
      "training_samples": "one",
      "agent_instruction": "..."
    },
    {
      "desc_level": "high",
      "task_type": "direct",
      "training_samples": "multiple",
      "agent_instruction": "..."
    },
    {
      "desc_level": "high",
      "task_type": "code",
      "training_samples": "multiple",
      "agent_instruction": "..."
    },
    {
      "desc_level": "none",
      "task_type": "direct",
      "training_samples": "none",
      "agent_instruction": "..."
    },
    {
      "desc_level": "none",
      "task_type": "code",
      "training_samples": "none",
      "agent_instruction": "..."
    },
    {
      "desc_level": "none",
      "task_type": "direct",
      "training_samples": "one",
      "agent_instruction": "..."
    },
    {
      "desc_level": "none",
      "task_type": "code",
      "training_samples": "one",
      "agent_instruction": "..."
    },
    {
      "desc_level": "none",
      "task_type": "direct",
      "training_samples": "multiple",
      "agent_instruction": "..."
    },
    {
      "desc_level": "none",
      "task_type": "code",
      "training_samples": "multiple",
      "agent_instruction": "..."
    }
  ],
  "observed_channels": [
    {"id": "Position", "label": "ball height", "unit": "m"},
    {"id": "Velocity", "label": "ball velocity", "unit": "m/s"},
    {"id": "Hard_Stop_f", "label": "contact impulse", "unit": "N*s"}
  ],
  "candidate_parameters": [
    "coefficient of restitution",
    "mass",
    "drag coefficient",
    "gravity acceleration"
  ],
  "download_link": "https://huggingface.co/datasets/eth-siplab/tracebench/tree/main/questions/BallDrop"
}
\end{DocCode}
The \code{agent_instruction} field is stored as plain text and must contain the fully rendered prompt shown to the agent.
It is rendered by \code{scripts/sync_hf_to_website_data.py} while synchronizing the GitHub Pages repository from Hugging Face.
The sync script reads \code{questions/<environment>/questions.json}, selects representative question entries for each website prompt condition, and renders \code{question_text} with \code{shared/prompts.py::render_tracebench_agent_prompt(...)} from the public TraceBench repository checkout \code{https://github.com/TommasoBendinelli/TraceBench.git}.

\phantomsection\label{sec:website-environment-data-json}
\begin{DocCode}
public/data/environments/BallDrop/data_1.json
{
  "run_id": "0f159f29c1004403a466c6bea03d7d02",
  "source": "TraceBench_questions/BallDrop/dataframes/0f159f29c1004403a466c6bea03d7d02.parquet",
  "columns": ["Position", "Velocity", "Hard_Stop_f", "time"],
  "intervention_time": 5,
  "intervention_parameter": "gravity",
  "rows": [
    {
      "Position": 10.882802,
      "Velocity": 5.839008,
      "Hard_Stop_f": 0.0,
      "time": 0.0025
    },
    {
      "Position": 10.897288,
      "Velocity": 5.794273,
      "Hard_Stop_f": 0.0,
      "time": 0.005
    },
    {
      "Position": 10.911662,
      "Velocity": 5.749537,
      "Hard_Stop_f": 0.0,
      "time": 0.0075
    }
  ]
}
\end{DocCode}
The \code{data_<n>.json} files are browser-ready JSON serializations of the Hugging Face files except \code{data_main_page.json} which is generated programmatically to match the GIF in the homepage (Only with ball height)
Specifically:
\begin{itemize}
  \item prompts, labels, and intervention metadata belong in the exported \code{questions.json}, \code{sample_manifest.json}, and \code{model_record.json} files.
  \item \agentProposedEdit{agent-remove-43f0272c-2980-5bbc-8fa4-e2f1cdbb3b37}{published sample paths now use questions/}{\code{TraceBench_questions/<environment>/dataframes/<run_id>.parquet} files.}{}
\end{itemize}

\appendix
\phantomsection\label{sec:website-prompt-appendix}
\section*{Appendix: Homepage Prompt Examples}

The homepage prompt panel is controlled by the \texttt{Task}, \texttt{Context}, and \texttt{Examples} buttons.
The \texttt{Noise} buttons change the plotted time series only; they do not change the prompt text.
Note, the homepage prompt is a simplified version of the real prompt, used to introduce the reader to the setup. 

\subsection*{Direct, High Context, No Examples}
\begin{DocCode}
The time series in the test_samples/ folder were generated by simulating a 1D vertical point-mass ball model under gravity.

For each simulation, either no parameter changes occur, or exactly one parameter among the allowed labels changes during the observed simulation interval.
If a parameter changes, it undergoes a single instantaneous step change at an unknown time during the observed interval.

Allowed labels:
["coefficient of restitution", "mass", "drag coefficient", "gravity acceleration", "no parameter change"]

Use "no parameter change" if there is no evidence in the data of a parameter change during the observed interval.

Task:
Create a file named results.json in the current working directory.
For each file in test_samples/, return the correct label.
The output must be valid JSON with exactly this structure:
{

  "<filename_1>": "<correct_label_filename_1>",
  "<filename_2>": "<correct_label_filename_2>"

}
\end{DocCode}

\subsection*{Direct, High Context, One Example}
\begin{DocCode}
The time series in the test_samples/ and train_samples/ folders were generated by simulating a 1D vertical point-mass ball model under gravity.

For each simulation, either no parameter changes occur, or exactly one parameter among the allowed labels changes during the observed simulation interval.
If a parameter changes, it undergoes a single instantaneous step change at an unknown time during the observed interval.

Allowed labels:
["coefficient of restitution", "mass", "drag coefficient", "gravity acceleration", "no parameter change"]

Use "no parameter change" if there is no evidence in the data of a parameter change during the observed interval.

Task:
Create a file named results.json in the current working directory.
For each file in test_samples/, return the correct label.
The output must be valid JSON with exactly this structure:
{

  "<filename_1>": "<correct_label_filename_1>",
  "<filename_2>": "<correct_label_filename_2>"

}

To help with this task, you can use the labeled train_samples/ directory containing one labeled example per class. 

\subsection*{Direct, High Context, Multiple Examples}
\begin{DocCode}
The time series in the test_samples/ and train_samples/ folders were generated by simulating a 1D vertical point-mass ball model under gravity.

For each simulation, either no parameter changes occur, or exactly one parameter among the allowed labels changes during the observed simulation interval.
If a parameter changes, it undergoes a single instantaneous step change at an unknown time during the observed interval.

Allowed labels:
["coefficient of restitution", "mass", "drag coefficient", "gravity acceleration", "no parameter change"]

Use "no parameter change" if there is no evidence in the data of a parameter change during the observed interval.

Task:
Create a file named results.json in the current working directory.
For each file in test_samples/, return the correct label.
The output must be valid JSON with exactly this structure:
{

  "<filename_1>": "<correct_label_filename_1>",
  "<filename_2>": "<correct_label_filename_2>"

}

To help with this task, you can use the labeled train_samples/ directory containing multiple labeled examples per class. 
\end{DocCode}

\subsection*{Code, High Context, No Examples}
\begin{DocCode}
Context:
The time series in the test_samples/ folder were generated by simulating a 1D vertical point-mass ball model under gravity.

For each simulation, either no parameter changes occur, or exactly one parameter among the allowed labels changes during the observed simulation interval.
If a parameter changes, it undergoes a single instantaneous step change at an unknown time during the observed interval.

Allowed labels:
["coefficient of restitution", "mass", "drag coefficient", "gravity acceleration", "no parameter change"]

Use "no parameter change" if there is no evidence in the data of a parameter change during the observed interval.

Task:
Create a Python script named rule.py in the current working directory.
The script must define exactly this function:
def predict(df) -> str:
The input df is a pandas DataFrame containing one sample with one column, 'ball height'.
The function will be called on samples inside test_samples/ and on additional held-out samples with the same schema and label set.
\end{DocCode}

\subsection*{Code, High Context, One Example}
\begin{DocCode}
Context:
The time series in the test_samples/ and train_samples/ folders were generated by simulating a 1D vertical point-mass ball model under gravity.

For each simulation, either no parameter changes occur, or exactly one parameter among the allowed labels changes during the observed simulation interval.
If a parameter changes, it undergoes a single instantaneous step change at an unknown time during the observed interval.

Allowed labels:
["coefficient of restitution", "mass", "drag coefficient", "gravity acceleration", "no parameter change"]

Use "no parameter change" if there is no evidence in the data of a parameter change during the observed interval.

Task:
Create a Python script named rule.py in the current working directory.
The script must define exactly this function:
def predict(df) -> str:
The input df is a pandas DataFrame containing one sample with one column, 'ball height'.
The function will be called on samples inside test_samples/ and on additional held-out samples with the same schema and label set.

To help with this task, you can use the labeled train_samples/ directory containing one labeled example per class while developing rule.py. 
\end{DocCode}

\subsection*{Code, High Context, Multiple Examples}
\begin{DocCode}
Context:
The time series in the test_samples/ and train_samples/ folders were generated by simulating a 1D vertical point-mass ball model under gravity.

For each simulation, either no parameter changes occur, or exactly one parameter among the allowed labels changes during the observed simulation interval.
If a parameter changes, it undergoes a single instantaneous step change at an unknown time during the observed interval.

Allowed labels:
["coefficient of restitution", "mass", "drag coefficient", "gravity acceleration", "no parameter change"]

Use "no parameter change" if there is no evidence in the data of a parameter change during the observed interval.

Task:
Create a Python script named rule.py in the current working directory.
The script must define exactly this function:
def predict(df) -> str:
The input df is a pandas DataFrame containing one sample with one column, 'ball height'.
The function will be called on samples inside test_samples/ and on additional held-out samples with the same schema and label set.

To help with this task, you can use the labeled train_samples/ directory containing multiple labeled examples per class while developing rule.py. 
\end{DocCode}

\subsection*{Direct, No Context, No Examples}
\begin{DocCode}
Context:
The time series in the test_samples/ folder were generated by a simulator of an unknown physical phenomenon.

For each simulation, either no parameter changes occur, or exactly one parameter corresponding to one of the allowed changes during the observed simulation interval.
If a parameter changes, it undergoes a single instantaneous step change at an unknown time during the observed interval.

["label_0", "label_1", "label_2", "label_3"] denote different parameter changes, while "label_4" denotes that no parameter changed.

Task:
Create a file named results.json in the current working directory.
For each file in test_samples/, return the correct label.
The output must be valid JSON with exactly this structure:
{

  "<filename_1>": "<correct_label_filename_1>",
  "<filename_2>": "<correct_label_filename_2>"

}
\end{DocCode}

\subsection*{Direct, No Context, One Example}
\begin{DocCode}
Context:
The time series in the test_samples/ and train_samples/ folders were generated by a simulator of an unknown physical phenomenon.

For each simulation, either no parameter changes occur, or exactly one parameter corresponding to one of the allowed changes during the observed simulation interval.
If a parameter changes, it undergoes a single instantaneous step change at an unknown time during the observed interval.

["label_0", "label_1", "label_2", "label_3"] denote different parameter changes, while "label_4" denotes that no parameter changed.

Task:
Create a file named results.json in the current working directory.
For each file in test_samples/, return the correct label.
The output must be valid JSON with exactly this structure:
{

  "<filename_1>": "<correct_label_filename_1>",
  "<filename_2>": "<correct_label_filename_2>"

}

To help with this task, you can use the labeled train_samples/ directory containing one labeled example per class. 
\end{DocCode}

\subsection*{Direct, No Context, Multiple Examples}
\begin{DocCode}
Context:
The time series in the test_samples/ and train_samples/ folders were generated by a simulator of an unknown physical phenomenon.

For each simulation, either no parameter changes occur, or exactly one parameter corresponding to one of the allowed changes during the observed simulation interval.
If a parameter changes, it undergoes a single instantaneous step change at an unknown time during the observed interval.

["label_0", "label_1", "label_2", "label_3"] denote different parameter changes, while "label_4" denotes that no parameter changed.

Task:
Create a file named results.json in the current working directory.
For each file in test_samples/, return the correct label.
The output must be valid JSON with exactly this structure:
{

  "<filename_1>": "<correct_label_filename_1>",
  "<filename_2>": "<correct_label_filename_2>"

}

To help with this task, you can use the labeled train_samples/ directory containing multiple labeled examples per class. 

\subsection*{Code, No Context, No Examples}
\begin{DocCode}
Context:
The time series in the test_samples/ folder were generated by a simulator of an unknown physical phenomenon.

For each simulation, either no parameter changes occur, or exactly one parameter corresponding to one of the allowed changes during the observed simulation interval.
If a parameter changes, it undergoes a single instantaneous step change at an unknown time during the observed interval.

["label_0", "label_1", "label_2", "label_3"] denote different parameter changes, while "label_4" denotes that no parameter changed.

Task:
Create a Python script named rule.py in the current working directory.
The script must define exactly this function:
def predict(df) -> str:
The input df is a pandas DataFrame containing one sample with one column, col1.
The function will be called on samples inside test_samples/ and on additional held-out samples with the same schema and label set.
\end{DocCode}

\subsection*{Code, No Context, One Example}
\begin{DocCode}
Context:
The time series in the test_samples/ and train_samples/ folders were generated by a simulator of an unknown physical phenomenon.

For each simulation, either no parameter changes occur, or exactly one parameter corresponding to one of the allowed changes during the observed simulation interval.
If a parameter changes, it undergoes a single instantaneous step change at an unknown time during the observed interval.

["label_0", "label_1", "label_2", "label_3"] denote different parameter changes, while "label_4" denotes that no parameter changed.

Task:
Create a Python script named rule.py in the current working directory.
The script must define exactly this function:
def predict(df) -> str:
The input df is a pandas DataFrame containing one sample with one column, col1.
The function will be called on samples inside test_samples/ and on additional held-out samples with the same schema and label set.

To help with this task, you can use the labeled train_samples/ directory containing one labeled example per class while developing rule.py. 
\end{DocCode}

\subsection*{Code, No Context, Multiple Examples}
\begin{DocCode}
Context:
The time series in the test_samples/ and train_samples/ folders were generated by a simulator of an unknown physical phenomenon.

For each simulation, either no parameter changes occur, or exactly one parameter corresponding to one of the allowed changes during the observed simulation interval.
If a parameter changes, it undergoes a single instantaneous step change at an unknown time during the observed interval.

["label_0", "label_1", "label_2", "label_3"] denote different parameter changes, while "label_4" denotes that no parameter changed.

Task:
Create a Python script named rule.py in the current working directory.
The script must define exactly this function:
def predict(df) -> str:
The input df is a pandas DataFrame containing one sample with one column, col1.
The function will be called on samples inside test_samples/ and on additional held-out samples with the same schema and label set.

To help with this task, you can use the labeled train_samples/ directory containing multiple labeled examples per class while developing rule.py.
\end{DocCode}

\phantomsection\label{sec:website-ball-drop-gif-prompt}
\section*{Appendix: BallDrop GIF Generation Prompt}

\begin{DocCode}
Create a deterministic animation of a 1D BallDrop simulator and its matching height-time trace that is generated by a python script.
 Output:
  - GIF: 4800×1800 px, 4 seconds, looping forever.
  - Also provide the full source script used to generate the files.
  - Final filenames:
    - ball-drop-bounce.gif
    - ball-drop-bounce-source.py
  - Implement with a 1600×600 logical canvas rendered at 3× scale to produce 4800×1800 output.

  Overall style:
  - White background.
  - Academic/research style.
  - Minimal black/gray linework.
  - One restrained blue accent only for the ball outline and plotted trace.
  - Orange only for the intervention marker.
  - No decorative elements, no mascot-like styling, no shadows, no gradients.
  - The GIF axis labels, tick labels, and intervention text must visually match the website
    Plotly example plot typography, including sans-serif font family, font weight, font size,
    text color, and label placement.

  Composition:
  - Split the 4800×1800 canvas into two unequal panels.
  - Left panel: compact simulator panel, approximately 18–22% of the canvas width.
  - Right panel: wide time-series plot panel, approximately 78–82% of the canvas width.
  - The left panel should feel intentionally compact because the ball only moves vertically.
  - The right plot should use the extra horizontal space to make the trace broad and easy to read.
  - Leave a subtle vertical divider or whitespace gap between the two panels.
  - Do not simply stretch an old 16:9 layout; redesign spacing for the wider 8:3 aspect ratio.

  Left panel:
  - Show a single outlined blue ball moving vertically.
  - The ball starts above the ground, not on the ground.
  - Show a thin black horizontal ground line near the bottom.
  - The ball should only move vertically; no horizontal movement.
  - Do not show any text labels in the simulator panel.
  - In particular, do not write “ground”.
  - Do not add a title above this panel.

  Right panel:
  - Show a simple time-series plot.
  - x-axis label: “time”
  - y-axis label: “ball height”
  - Use thin black axes.
  - The plotted blue trace should reveal over time.
  - The plotted trace must match the ball’s vertical position exactly.
  - Do not add a plot title above this panel.

  Physics:
  - Simulate vertical motion under gravity.
  - Use deterministic trajectory generation; no randomness.
  - The ball starts at a positive height above the ground.
  - Use a bounce-height coefficient, not a velocity coefficient.
  - The first 3 rebounds must use e = 0.9, so each rebound peak height is 90% of the previous peak height.
  - Immediately after the third ground impact, mark the intervention time on the plot and switch the bounce-height coefficient
  to e = 0.2.
  - Show a vertical orange intervention marker on the plot at the intervention time.
  - Label the marker with:
    - “intervention”
    - “e: 0.9 -> 0.2”
  - The intervention label should be large enough to read at webpage display size.
  - After the intervention, continue simulating repeated bounces using e = 0.2.
  - Do not stop after only one post-intervention bounce.
  - Each post-intervention bounce should reach 20% of the previous bounce height, producing a rapidly decaying sequence of small
  bounces.
  - The post-intervention motion should not appear as one small bounce followed by an immediate snap to zero.
  - Only set the ball to rest on the ground after the bounce height becomes visually negligible, for example less than 0.005 of
  the plot height.
  - Then keep the ball resting on the ground for more than 1 second before the animation loops.
  - The total rendered animation must still be exactly 4 seconds.

  Animation behavior:
  - The ball motion and the trace reveal should be synchronized frame-by-frame.
  - The final portion may show the ball resting on the ground, but avoid making the animation feel unnecessarily static.
  - Use a fixed frame rate and deterministic sampling so repeated runs produce the same result.
  - The ball should never move horizontally.
  - The plotted trace and the ball’s vertical position must come from the same trajectory data.

  Implementation requirements:
  - Generate the trajectory numerically or analytically in code.
  - Generate all frames deterministically from the same trajectory data.
  - Do not manually draw inconsistent ball and plot positions.
  - Export:
    - ball-drop-bounce.gif
    - ball-drop-bounce-source.py
\end{DocCode}


\end{document}
